\documentclass[a4paper,11pt]{article}
\usepackage[utf8]{inputenc}
\usepackage[english]{babel}
\usepackage{geometry}
\usepackage{titlesec}
\usepackage{enumitem}
\usepackage{hyperref}
\usepackage{graphicx}

\geometry{margin=2.5cm}

\title{
    \vspace{-2cm}
    \textbf{Software Development Fundamentals (FSOFT)}\\
    \large Technical Report - Iteration 01 (M1)\\
    \Large \textbf{Supermarket Inventory and Sales Management System}
}

\author{
    \begin{tabular}{c}
    Team: [1DB\_3] \\
    Members: \\
    Tomás Bastos 1251506, Igor Almeida 1251054 \\
    Guilherme Asevedo 1251002, Rui Silva 1251443 \\
    ISEP - DEI
    \end{tabular}
}
\date{March 2026}

\begin{document}

\maketitle

\begin{abstract}
This technical report describes the first iteration (M1) of the project to develop an inventory and sales management application for a supermarket. It focuses on requirements analysis, identification of actors, and definition of use cases, following Software Engineering principles and Object-Oriented Programming in C++.
\end{abstract}

\section{Introduction}
The objective of this project is to develop a robust C++ application to manage the daily operations of a supermarket, including stock control, sales management, and customer loyalty. This technical report details the first iteration (M1), which corresponds to the \textit{Requirements Engineering} and \textit{Analysis} phases of the Software Development Life Cycle (SDLC). The primary goal is to establish a solid functional foundation through Object-Oriented Analysis (OOA), ensuring that the software alignment with stakeholder needs.

\section{Software Development Process}
Following the FSOFT course guidelines, this project adopts an \textbf{Iterative and Incremental} development model. This approach allows for continuous refinement of requirements and design across multiple cycles (Iterations 01 to 03).

\subsection{Methodology and Tools}
\begin{itemize}
    \item \textbf{Process:} Object-Oriented Analysis (OOA) was employed to identify domain concepts and system behaviors.
    \item \textbf{Version Control:} Git is used as the Version Control System (VCS), with a remote repository hosted on GitHub to facilitate collaboration and maintain a historical record of changes, adhering to best practices such as descriptive commit messages.
    \item \textbf{Modelling:} Unified Modeling Language (UML) is used to represent the system's structural and behavioral aspects.
\end{itemize}

\section{Requirements Analysis (OOA)}
In this phase, we acquired domain knowledge and reasoned about the supermarket business logic to define the system scope.

\subsection{System Actors}
We have identified the following actors (primary and secondary) interacting with the system boundary:
\begin{itemize}
    \item \textbf{Administrator (Admin):} A primary actor with full privileges. Responsible for managing the product catalog, user accounts, promotions, and auditing the system through logs and reports.
    \item \textbf{Cashier:} A primary actor responsible for the point-of-sale (POS) operations, including recording sales, managing customer loyalty data, and issuing receipts.
    \item \textbf{System:} A secondary (temporal) actor that automatically monitors background events, such as low stock levels and product expiration dates.
\end{itemize}

\subsection{Use Cases}
The functional requirements are mapped into Use Cases, representing the goals of the actors:
\begin{description}
    \item[UC01: Manage Products] (Admin) Create, update, delete, and list products (Code, Category, Price, Expiry).
    \item[UC02: Manage Users] (Admin) Manage employee access levels and credentials.
    \item[UC03: Manage Promotions] (Admin) Define discount rules applied during the sale process.
    \item[UC04: Consult Reports] (Admin) Analyze aggregated sales data and performance metrics.
    \item[UC05: Consult Logs] (Admin) Audit system changes for security and accountability.
    \item[UC06: Register Sale] (Cashier) Process a new transaction, identifying products and applying promotions.
    \item[UC07: Issue Receipt] (Cashier) Generate a formal proof of purchase for the customer.
    \item[UC08: Manage Customers] (Cashier) Register customers via NIF for loyalty points or tax identification.
\end{description}

\section{Domain Model Discovery}
As part of the OOA primary tasks, we identified the following core domain objects that will be further detailed in the Design phase (Iteration 02):
\begin{itemize}
    \item \textbf{Product:} Represents items in stock with attributes like \texttt{ID}, \texttt{description}, \texttt{stockQuantity}, and \texttt{expiryDate}.
    \item \textbf{Sale:} Represents a transaction, composed of multiple \texttt{SaleItems}.
    \item \textbf{User:} Represents the system's users (Admin/Cashier) with \texttt{username} and \texttt{hashedPassword}.
    \item \textbf{Promotion:} Encapsulates discount logic.
\end{itemize}

\section{System Requirements}

\subsection{Functional Requirements (FR)}
Derived from the use cases, the system must fulfill:
\begin{enumerate}[label=FR\arabic*:]
    \item The system shall provide persistent storage for product data (category, price, expiration).
    \item The system shall trigger alerts for low stock (below a defined threshold) or near-expiry items.
    \item The system shall calculate the subtotal and final total of a sale in real-time.
    \item Data persistence must be implemented using plain text files (\texttt{.txt} or \texttt{.csv}).
    \item Every price modification must be recorded in an immutable audit log.
\end{enumerate}

\subsection{Non-Functional Requirements (NFR)}
Reflecting the technical constraints of the project:
\begin{enumerate}[label=NFR\arabic*:]
    \item \textbf{Language:} Implemented exclusively in C++, utilizing modern standards (C++11/14) and safe memory management (avoiding raw pointers where possible, as per T2).
    \item \textbf{Interface:} A text-based Command Line Interface (CLI) organized with hierarchical menus.
    \item \textbf{Platform:} Compatibility with Linux/Unix environments (targeted at ISEP laboratory environments).
    \item \textbf{Security:} Role-based access control (RBAC) requiring authentication for all sensitive operations.
\end{enumerate}

\section{UML Artifacts}
The Use Case Diagram (UCD) serves as the primary artifact for this iteration, visualizing the relationship between actors and system functionalities.
\textit{[Note: The Use Case Diagram image should be placed here]}

\section{Conclusion}
In this first iteration, we established the functional foundation of the Supermarket Management System. By applying OOA principles, we clearly defined the actors and requirements. This analysis provides the necessary blueprint for Iteration 02, where we will move into Object-Oriented Design (OOD) and the initial C++ implementation.

\end{document}
